% !TeX root = closed-systems-from-comparison-completeness.tex
% ============================================================
\chapter*{Logical Status and Dependency Ledger}
\phantomsection
\addcontentsline{toc}{chapter}{Logical Status and Dependency Ledger}
\label{ch:logical-status}
% ============================================================
The monograph distinguishes definitions, intrinsic consequences,
realization-level statements, boundary data, and open reconstruction
problems.  This distinction is part of the formal content of the book:
a condition is not treated as primitive merely because it appears in a
later theorem header.

\section*{Status classes}
\begin{description}[leftmargin=3.5cm,style=nextline]
\item[Primitive definition]
A stipulation fixing the object of study.  The central primitive is the
comparison world \((U,\mathcal C)\), with intrinsicality built into what
counts as an admissible comparison predicate.
\item[Intrinsic theorem]
A result proved from comparison data, closed-world admissibility, quotient
descent, and the transport hierarchy.
\item[Scoped setup]
Local objects, indices, connectedness assumptions, or finite diagrams
introduced to state a theorem.  These do not function as new
foundational primitives.
\item[Realization-level theorem]
A theorem proved after a smooth, Hilbert, phase, observer-side, or
quantitative realization has been introduced.
\item[Boundary condition]
Data explicitly retained at an interface because it is not derived in
the intrinsic stack at that point.
\item[Open bridge]
A named reconstruction or quantitative step not claimed as a completed
intrinsic theorem.
\end{description}

\section*{Principal dependency chain}
\begin{enumerate}[label=\textup{(\arabic*)}]
\item \textbf{Primitive comparison world.}
The primitive object is \((U,\mathcal C)\), not a manifold, field bundle,
Hilbert space, or causal order.  See
\cref{def:intro-comparison-world,ch:foundation}.

\item \textbf{Finite-to-global admissibility.}
The refinement tower supplies the finite witness and compactness boundary
used by later constructions.  See \cref{ch:foundation}.

\item \textbf{Comparison completeness and quotient semantics.}
Under intrinsicality and local distinguishability, rectangular
completeness is derived; quotient semantics then determine the physical
content that descends through intrinsic redundancy.  See
\cref{ch:bottom,ch:rotation,ch:closure}.

\item \textbf{Two-locus classification.}
Every enrichment beyond quotient semantics enters through representative
choice, transport, or both.  See \cref{ch:locus,ch:lift}.

\item \textbf{Transport obstruction.}
Finite triangle failure, global loop obstruction, and the bridge between
them identify the first stable visible obstruction layer.  See
\cref{ch:triangles,ch:descent_obstruction,ch:bridge}.

\item \textbf{Observer reduction and spacetime interleaving.}
The obstruction stack supports internal observer reduction and a
two-parameter inverse-limit reading of spacetime interleaving.  See
\cref{ch:flow,ch:stitching}.

\item \textbf{Intrinsic--smooth interface.}
Smooth geometry enters as a realized reading.  Faithfulness on the forced
degree--\(2\) carrier is analyzed by the torsion condition \textup{(T)}
and detectability axiom \textup{(D)}, with ghost counter-models marking
necessity and a canonical realization discharging the condition in the
specified regime.  See \cref{ch:christoffel,ch:interface,ch:connection}.

\item \textbf{Macroscopic, Hilbert, electromagnetic, and scalar sectors.}
Causal scaling, Hilbert structure, Einstein compatibility,
connection-first reconstruction, phase curvature, charge typing, and
scalar-sector closure are downstream realization-level consequences of
the transport and interface stack.  See
\cref{ch:causality,ch:hilbert,ch:einstein,ch:connection,ch:em,ch:matter,ch:matter-numbers,ch:numerical-closure,ch:matter-realization}.

\item \textbf{State-side reconstruction extension.}
After the realized stack is in place, the raw limit comparison algebra is
shown to exceed any realized smooth reading.  Under scale coherence
\textup{(U)}, the regular sector and its derivations are identified.  See
\cref{ch:dynamics-program}.

\item \textbf{Quantitative and application-level completion.}
Remaining global quantitative data are recorded in the conditional
ledger, and the scalar channel is applied to a concrete confinement
realization.  See \cref{app:conditional-completion,app:fusion}.
\end{enumerate}

\section*{Claims not made}
The monograph does not claim that a smooth manifold, Hilbert space,
field bundle, or observer-side measurement scale is primitive.  It also
does not claim that every quantitative physical parameter is derived in
the intrinsic stack.  When such data enter, the text marks them as
realization data, boundary conditions, or open reconstruction bridges.

\section*{Use in later chapters}
When a later theorem states a smoothness, faithfulness, ghost-free,
scale-coherence, or observer-readout condition, that condition should be
read according to this ledger.  It is part of the theorem's domain of
validity, not an unmarked return to background-first modeling.
